\section{LaTeX Integration}
\label{sec:latex-integration}

Triton diagrams can be embedded in \LaTeX{} documents as \emph{vector} PDF
figures. The integration lives in its own satellite package, \texttt{latex/}
(\texttt{@triton/latex}), and is purely additive: it changes no core behaviour and
adds \textbf{zero} dependencies to the \texttt{triton} package.

\subsection{The format gap}

\LaTeX{}'s \texttt{\textbackslash includegraphics} ingests PDF, PNG, or JPG ---
never SVG --- on every engine (pdf\LaTeX{}, Xe\LaTeX{}, Lua\LaTeX{}). Triton emits
SVG (\S\ref{sec:scene}: \texttt{renderSync} $\rightarrow$ \texttt{renderSVG}). The
integration therefore converts each diagram to a PDF \emph{ahead of time}, and the
document includes the committed PDF. This ``precompile and commit'' model is the
only one that also works on Overleaf, where no Node or Triton toolchain is
available at compile time.

\subsection{SVG \texorpdfstring{$\rightarrow$}{->} vector PDF, pure-JS}

The converter is pure JavaScript with \textbf{no system binaries} (no Inkscape, no
\texttt{rsvg-convert}, no headless browser). It feeds Triton's SVG through
\texttt{pdfkit} + \texttt{svg-to-pdfkit}, producing a single-page PDF whose page
box equals the diagram's \texttt{viewBox}. The output is genuinely vector:

\begin{itemize}
  \item shapes become PDF path operators (\texttt{re}, \texttt{l}, \texttt{c});
  \item the \texttt{Scene}'s five element types and the arrowhead
        \texttt{<marker>} definitions all replay faithfully;
  \item text becomes real vector glyphs in the embedded base-14 fonts
        (Helvetica / Times / Courier) --- so there is \emph{no font drift}
        between author, CI, and Overleaf, and no external font file to ship.
\end{itemize}

The SMIL animation overlays (\texttt{march}, \texttt{particle}) are dropped: a PDF
is static, and the first frame is the correct print representation.

\subsection{The CLI}

The package ships a \texttt{triton-latex} command (esbuild-bundled from
\texttt{latex/src/cli.ts}, which imports the compiler by relative path from
\texttt{../src} --- the same satellite pattern as the VS Code extension):

\begin{lstlisting}[language=ts, basicstyle=\small\ttfamily]
triton-latex render <input.triton|.mmd> -o <out.pdf|.svg> [--theme N] [--scale S]
triton-latex render-dir <srcDir>        -o <outDir>        [--theme N] [--scale S]
\end{lstlisting}

\texttt{render} converts one source to vector PDF (or passes SVG through);
\texttt{render-dir} batch-renders every \texttt{*.triton}/\texttt{*.mmd} in a
directory to \texttt{<name>.pdf}. Both reuse the core \texttt{renderSync()} path
and surface its \texttt{Result} as a clean exit code.

\subsection{The \texttt{triton.sty} package}

A thin \LaTeX{} package wraps \texttt{\textbackslash includegraphics}, resolving a
diagram by name against a configurable directory. It depends only on
\texttt{graphicx}, so it is engine-agnostic.

\begin{lstlisting}[language=ts]
\usepackage{triton}
\tritondir{figures}                 % where the rendered <name>.pdf live
\triton{flowchart}                  % includegraphics figures/flowchart.pdf
\tritonfig[width=0.6\linewidth]{avl}{An AVL tree rendered by Triton.}
\end{lstlisting}

\subsection{Workflow and portability}

The authoring loop mirrors the design doc's own \texttt{pnpm figures} /
\texttt{\textbackslash ourfig} convention: render sources to PDF, commit them,
include by name, and regenerate when a diagram changes. Because the assets are
committed vector PDFs, the same project compiles unchanged on macOS, Linux,
Windows, and Overleaf --- the engine axis is irrelevant, since the asset is
already a PDF by the time \texttt{\textbackslash includegraphics} runs.

\begin{center}
\begin{tabular}{ll}
\toprule
Path & Contents \\
\midrule
\texttt{latex/src/cli.ts}   & the \texttt{triton-latex} CLI (render / render-dir) \\
\texttt{latex/src/pdf.ts}   & SVG $\rightarrow$ vector PDF (pdfkit + svg-to-pdfkit) \\
\texttt{latex/triton.sty}   & the \texttt{\textbackslash triton} / \texttt{\textbackslash tritonfig} macros \\
\texttt{latex/esbuild.mjs}  & bundles the CLI + compiler graph from \texttt{src/} \\
\texttt{latex/examples/}    & \texttt{diagrams/} sources, committed \texttt{figures/} PDFs, \texttt{demo.tex} \\
\bottomrule
\end{tabular}
\end{center}
